\section{Pass protocol and composition}
\label{sec:composition}\label{sec:staging}
A macro invocation carries a complete program and an explicit queue of remaining
pass paths, each with argument tokens read only by that pass.
\begin{mathpar}
\Rule{ST-Entry}{\pi=\mathsf{resolve}(configuration,defaults)\quad\pi\ne[]}
 {\code{mini_linq::compile!}\{configuration;\ T\}
  \longrightarrow\mathsf{start}(\pi,T)}

\Rule{ST-Next}{\operatorname{parse}_L(T)=s\quad\mathsf{WF}_L(s)\\
 f(s,args)=t\quad\mathsf{WF}_M(t)\quad M\ne Rust\quad\pi=k(args')::\pi'}
 {\mathsf{pass}_f(args,\pi,T)
  \longrightarrow k!\{args';\ \pi';\ \mathsf{quote}(t)\}}

\Rule{ST-Context}{\mathsf{pass}_f(args,\pi,\mathsf{quote}(s))\longrightarrow T'}
 {C[\mathsf{pass}_f(args,\pi,\mathsf{quote}(s))]\longrightarrow C[T']}

\Rule{ST-Sequence}{q=\mathsf{sequence}(s)\quad\mathsf{check}(s,q)\\
 \mathsf{checkGrading}(q,host)\quad e=\mathsf{translatePull}(q)\quad\pi=[]}
 {\mathsf{sequence};\mathsf{sequential};\mathsf{translatePull}(s)
  \longrightarrow\mathsf{quote}_{Rust}(e)}
\end{mathpar}
\code{cq_stage::pass} checks source and target contracts; \code{cq_stage::rewrite}
also runs the language's \code{check_rewrite}. Correspondence checks belong to the
pass. Failures emit \code{compile_error!}; rustc then checks types, ownership and
trait obligations of the generated code.
\begin{lstlisting}
pipeline {
    surface = [Path(tokens), ...];
    plan = Path(tokens); ix = Path(tokens); lower = Path(tokens);
    physical = [Path(tokens), ...]; select = Path(tokens);
    structures = [Path(tokens), ...];
    represent = Path(tokens); stage = Path(tokens);
    sequence = Path(tokens); decompose = [Path(tokens), ...];
    grade = Path(tokens); print = Path(tokens);
}
\end{lstlisting}
The order is surface, plan, ix, lower, Physical rewrites, select, Structures
rewrites, represent, stage, sequence, decompose, grade, print. The shipped defaults
are \code{plan}, \code{index_demands}, \code{lower}, \code{select},
\code{represent}, \code{stage}, \code{sequence}, \code{sequential}, and
\code{translate_pull}. Thus ST-Sequence is the default CPU path.
Decompose is an optional list of Sequence rewrites.
An explicit Structures terminal such as \code{own_push::push}, with no
authored sequence, grade, or decompose phases, bypasses the default Sequence
phases. Custom Sequence terminal paths require explicit sequence and grade
phases; terminal compatibility is checked by the pass protocol.
\code{translate_pull} emits both ordinary and profile-counter variants.
The workload corpus and fuzz configuration use this target by default.
An authored fact-free program may enter at any compatible macro with
\code{arguments \{..\}; pipeline \{..\}; <program>}. Passes receive no
configuration beyond their arguments: every choice is in the program text.

\paragraph{Context envelopes at invocation boundaries.}
The rules above retain the fact-free protocol. When a source already records
discharged host facts or a consumer fact (Section~\ref{sec:context}), its
invocation must instead carry an explicit \code{context \{..\};} envelope
between the optional arguments and the pipeline. The envelope must equal
the recorded source facts in both directions: it may neither omit a recorded
fact nor add an unrecorded one. Locally certified host facts are not envelope
facts. A lower/select argument packet may introduce facts into a fact-free
source; naming facts in such a packet never substitutes for the required
envelope on an already-factful source. The enclosing compiler that emits a
hole invocation must write that envelope and prove its facts.

Each invocation consumes and checks its own envelope. In \textsc{ST-Next},
the emitted invocation derives a fresh envelope from the recorded facts of
$t$, its next source, rather than blindly forwarding the consumed envelope;
it omits the envelope if those facts are empty. Physical facts carried by
containment into Structures and Sequence preserve the envelope when those
facts are unchanged. Ordinary passes and rewrites may preserve or drop facts,
not introduce them; only the lower/select packet boundary authorizes
introduction. Missing or inexact envelopes and unauthorized introductions
are rejected as \textsc{CTX-Entry} failures.

Optional features are standalone workspaces under \code{features/} that depend
on the packaged core only: planners, rewrites, surface extensions (negation,
aggregation, correlated subqueries, recursion as Rust contexts around checked
core invocations), a push printer, a stacked IR, a diverged RA surface,
composition, inspection and workloads. \code{tools/check_external.py} packages
the core and builds and tests each feature outside the repository, and rejects a
fixture that only builds in-tree.

\begin{theorem}[Conditional composition]
\label{thm:composition}
If $f:L\to M$ and $g:M\to N$ terminate on their domains, $f$ produces only inputs
$g$ admits and establishes $g$'s representation preconditions, and both preserve
a common observation under compatible closing contexts, then so does $g\circ f$.
\end{theorem}
\begin{proof}Compose the representation relations and observation equalities;
domain inclusion makes $g$ defined; context adequacy lifts the argument to macro
invocations.\end{proof}

\paragraph{Set aggregate contexts.}
The aggregate context records its row-stream canonicality proof, owned head
columns, constant-group proof and \code{cq_storage::HashState}. The default
Sequence translator preserves the checked canonical-return boundary; the extension folds
that stream directly without adding another distinct set. Unknown terminals
may repeat rows and retain an explicit fast-hashed distinct boundary.
\begin{mathpar}
\Rule{EXT-SetFold}{\mathsf{checkedCanonicalTerminal}(\pi)\quad
 Q:\bot+\mathcal P(Row)}
 {\mathsf{aggregate}(Q)=\mathsf{fold}(Q)}

\Rule{EXT-ConstantGroup}{\forall k\in K.\ \mathsf{literalBound}(k)}
 {\mathsf{groupFold}_K(Q)=\mathsf{optionalAccumulator}(Q)}
\end{mathpar}
Constant grouping emits \code{SingleGroup}, which stores zero or one key and
measure tuple without hashing. Empty input still has no group. Computed measures
are borrowed for cloning at the fold, rather than materialized as another CQ
input. General groups use the recorded fast hash state, including their reusable
answer storage. None of these facts implies input-row or primary-key uniqueness:
the core's necessary result deduplication remains.
